% Verbatim copy of /Users/alingalatan/workspaces/grothe/handoff/manuscript/legacy/ch47-cylinder-lifting-section.tex (file dated 2026-10-25). Only this line was added.
\subsection{A finite cylinder lifting in arbitrary dimension}
\label{sec:uniform-complement}
Let $J\in\mathcal H_{2,1}$ satisfy $\mathbb E[J\mid Z,X]=0$ and let the fresh buffer be independent of every coordinate of $J$. The following algebra explains why the core cell majorants can control arbitrary additional dimensions. No independence of the core and complement is used.
\paragraph{Step 1: retain both signs.}
On a cell measurable in $(Z,X)$, let $q_{j,\pm}$ be the incoming
anchor second moments and put $F_\pm=L\pm Q$. Concavity of the Rice
mean gives
\[
\begin{split}
 \mathbb E_j\rho_g(|F_\pm\pm J|^2)
 \le{}&\rho_g(q_{j,\pm})+\frac{K_g(q_{j,\pm})}{2}
 \bigl(\mathbb E_j|F_\pm|^2-q_{j,\pm}+\mathbb E_j|J|^2\bigr).
\end{split}
\]
Conditional centering eliminates the cross term. Both signs are needed:
$Q+J$ need not have the one-driver sign symmetry. Their extra contribution
is $\mathbb E[W_g|J|^2]$, with the cellwise weight
\[
 W_g=\tfrac12\bigl(K_g(q_{j,+})+K_g(q_{j,-})\bigr).
\]
Set $W_g=1$ outside the bounded cell domain. The original inequality
$2|w|\le1+|w|^2$ gives this same extra tail contribution; conditional
centering eliminates tail cross terms as well. Orthogonality supplies
an additional penalty $-p\|J\|_2^2$. Thus every original majorant
$M(Q,g)$ satisfies
\[
 \Lambda(Q+J+gU)\le M(Q,g)+\mathbb E[(W_g-p)|J|^2].
\]
Moreover
\[
 K_g(u)=\frac1{\sqrt\pi}\int_0^{1/g^2}
 e^{-ut}t^{-1/2}\sqrt{1-g^2t}\,dt\le u^{-1/2}=K_0(u),
\]
so $W_g\le W_0$ for every buffer.

\paragraph{Step 2: five matrices suffice in every dimension.}
Expand in the normalized complex Hermite basis of the additional
coordinates $V$:
\[
 J=\sum_{\gamma,\delta}j_{\gamma\delta}(Z,X)
 \mathsf H_{\gamma\delta}(V),\qquad
 j_{\gamma\delta}\in
 \mathcal H_{2-|\gamma|,1-|\delta|}(\mathbb C^2).
\]
The constant remote label is zero by conditional centering. Orthogonality
and the fact that $W_0$ depends only on $(Z,X)$ give
\[
 \mathbb E[W_0|J|^2]
 =\sum_{\gamma,\delta}\mathbb E[W_0|j_{\gamma\delta}|^2],
 \qquad
 \|J\|_2^2=\sum_{\gamma,\delta}\|j_{\gamma\delta}\|_2^2.
\]
The possible coefficient spaces are precisely
\[
 \mathcal H_{1,1},\quad\mathcal H_{2,0},\quad
 \mathcal H_{0,1},\quad\mathcal H_{1,0},\quad\mathcal H_{0,0},
\]
of dimensions $4,3,2,2,1$. Arbitrarily many remote labels repeat these
same matrices, rather than creating a dimension-dependent estimate.

\paragraph{Step 3: explicit block entries.}
Remove the common phase and write $Z=r$, $X=x+iy$, $s=r^2$,
$t=x^2+y^2$. For a real anchor the paired weight is even in $x$ and $y$ separately; the next section retains the complex entries for nonreal anchors.
Write $\mathbb E_W f=\mathbb E[W_0f]$. In the orthonormal basis
$(s-1,Z\overline X,X\overline Z,t-1)$, the $\mathcal H_{1,1}$ matrix
splits into
\[
 \begin{pmatrix}
 \mathbb E_W(s-1)^2&\mathbb E_W(s-1)(t-1)\\
 \mathbb E_W(s-1)(t-1)&\mathbb E_W(t-1)^2
 \end{pmatrix},\qquad
 \begin{pmatrix}
 \mathbb E_Wst&\mathbb E_Ws(x^2-y^2)\\
 \mathbb E_Ws(x^2-y^2)&\mathbb E_Wst
 \end{pmatrix}.
\]
In the basis $(Z^2/\sqrt2,ZX,X^2/\sqrt2)$, the $\mathcal H_{2,0}$
matrix splits into
\[
 \begin{pmatrix}
 \mathbb E_Ws^2/2&\mathbb E_Ws(x^2-y^2)/2\\
 \mathbb E_Ws(x^2-y^2)/2&\mathbb E_Wt^2/2
 \end{pmatrix}
\]
and the singleton $\mathbb E_Wst$. The other matrices are
$\operatorname{diag}(\mathbb E_Ws,\mathbb E_Wt)$, its conjugate,
and the singleton $\mathbb E_W1$.

Each matrix is evaluated as
\[
 I+\sum_{j\subset D}p_j(w_j-1)\mathbb E_j[\Phi\Phi^*].
\]
This incorporates the unbounded complement exactly, since $W_0=1$
there. It only needs conditional polynomial moments of degree at most
four, in addition to the original anchor moments of degree six.

